\documentclass{situationist-preprint}

\title{Akrasia and Procrastination}
\msid{ACT-04-b}
\msfamily{ACT-04}
\msdiscipline{Action, Agency and Free Will}
\mstype{Consequence}
\mskeywords{weakness of will; akrasia; procrastination; motivational strength; temporal discounting}

\begin{document}
\maketitle

\begin{abstract}
If strict akrasia is possible, as ACT-04 argues, what follows for procrastination? We derive three consequences. First, procrastination, understood as voluntary delay against one's better judgment, is possible, and its possibility requires no premise beyond those that make akrasia possible. Second, procrastination is a species of akrasia but not the only one, since akrasia need not involve delay. Third, the account of akrasia explains the characteristic temporal shape of procrastination: the agent's judgment is constant while the motivational strength of the competing options changes with their proximity. The procrastinator is motivated by her judgment, just not enough, and not enough at the moment of action.
\end{abstract}

\section{From akrasia to procrastination}

ACT-04-a defines strict akrasia as intentionally and freely performing an option other than the one one's better judgment favors, while holding that judgment (Definition~2.2), and proves that it is possible (Theorem~3.2). The key premise is that a better judgment motivates the agent who holds it to some degree, but that its motivational strength is not fixed by its content (Principle~P1). The akratic agent is motivated by her judgment, just not enough. Procrastination is commonly defined as voluntarily delaying an intended course of action despite expecting to be worse off for the delay \citep{Steel2007}; ACT-13-a makes the expectation explicit as a better judgment held at the time of delay.

\begin{proposition}\label{p:poss}
Procrastination is possible.
\end{proposition}

\begin{proof}
Let the options at $t$ be beginning a task and delaying it, and let the agent judge at $t$ that beginning is better, all things considered. By Proposition~3.1 of ACT-04-a, she may be more strongly moved to delay, since the strength of her judgment's motivation is not fixed by its content, and the strength of the desire to avoid the task's present costs is fixed by other facts. By the argument of Theorem~3.2 there, she may then delay intentionally and freely while holding her judgment. If she had earlier intended to perform the task, this is procrastination.
\end{proof}

\begin{proposition}\label{p:species}
Every case of procrastination is a case of akrasia, but not conversely.
\end{proposition}

\begin{proof}
The first claim is Lemma~3.1 of ACT-13-a. For the second, an agent who judges that she should decline a second glass of wine and accepts it acts akratically, but does not delay any task she intends to perform; her akrasia concerns what she does, not when she does it.
\end{proof}

\section{The temporal shape}

The account of ACT-04-a does more than permit procrastination; it explains its shape. ACT-04-a observes that the attractiveness of an option varies with its temporal proximity in a way that its assessed value does not. A task's costs fall when it is begun and its benefits later. If motivational strength is discounted steeply for short delays \citep{Ainslie1992}, the pull of avoiding the costs rises sharply as the time to begin approaches, while the agent's judgment of the value of beginning, which does not track proximity as such, remains fixed. Before the time comes, judgment and motivation agree and the agent forms the intention; at the time, they diverge, and the agent delays. This is the mechanism that ACT-13-a identifies in its Lemma~3.2. It requires the separation of evaluative ranking from motivational strength that ACT-04-a defends, and it is unavailable to a view that identifies judging an option better with wanting it more. On such a view the procrastinator's judgment would have to reverse when her motivation does, and she would not procrastinate but change her mind.

\begin{objection}[Change of mind]
The procrastinator does change her mind: at the moment of delay she judges that delaying is better, and later regrets it.
\end{objection}

\begin{reply}
Some delayers do, and they do not procrastinate; they postpone, and their later regret is a change of view. The procrastinator is characterized by holding her judgment through the delay, as her own reports confirm: she says, at the time, that she should be working. To redescribe every such case as a fleeting change of judgment is to deny the phenomenon, and the denial rests on the identification of judgment with motivation that ACT-04-a rejects.
\end{reply}

\section{Conclusion}

Since akrasia is possible, so is procrastination, which is akrasia whose object is delay and whose source is the changing motivational strength of options as they approach. The account of ACT-04 thus supplies both the possibility and the mechanism of procrastination, and with them the ground of ACT-13's conclusion that procrastination is irrational.

\begin{thebibliography}{9}
\bibitem[Ainslie(1992)]{Ainslie1992} Ainslie, G. (1992). \emph{Picoeconomics: The Strategic Interaction of Successive Motivational States Within the Person}. Cambridge University Press.
\bibitem[Steel(2007)]{Steel2007} Steel, P. (2007). The nature of procrastination: a meta-analytic and theoretical review of quintessential self-regulatory failure. \emph{Psychological Bulletin}, 133(1), 65--94.
\end{thebibliography}

\end{document}
